%\documentclass[icelandic]{beamer}

\documentclass[icelandic,a4paper,12pt]{article}
\usepackage{beamerarticle}

\mode<presentation>
{
  \usetheme{boxes}
  % með efnisyfirliti: Szeged, Frankfurt 
  % án efnisyfirlits: Pittsburgh
  % áhugavert: CambridgeUS, Boadilla
  %\setbeamercovered{transparent} %gegnsætt
  \setbeamercovered{invisible}
  }

\usepackage[english,icelandic]{babel}
\usepackage[utf8]{inputenc}
\usepackage{t1enc}
\selectlanguage{icelandic}
\usepackage{graphicx}
\usepackage{amsmath}
\usepackage{amssymb}
\usepackage{mathrsfs}
% \newcommand{\C}{{\mathbb  C}}
% \newcommand{\Z}{{\mathbb Z}}
% \newcommand{\R}{{\mathbb  R}}
% \newcommand{\N}{{\mathbb  N}}
% \newcommand{\Q}{{\mathbb Q}}
\renewcommand{\phi}{\varphi}
\renewcommand{\epsilon}{\varepsilon}

%\usepackage{pgfpages}
% \pgfpagesuselayout{2 on 1}[a4paper,border shrink=5mm]

\def\lecturename{Stærðfræðigreining IB}
\title{\insertlecture}
\author{Benedikt Steinar Magnússon, \href{mailto:bsm@hi.is}{bsm@hi.is}}
\institute
{
  Verkfræði- og náttúruvísindasvið\\
  Háskóli Íslands
}
\subtitle{Stærðfræðigreining IB}
%\subject{\lecturename}

\mode<article>
{
	\usepackage[colorlinks=false,
	pdfauthor={Benedikt Steinar Magnusson},
	pdftitle={IB: Namsefni
	}]{hyperref}
  %\usepackage{times}
  %\usepackage{mathptmx}
  \usepackage[left=1.5cm,right=4cm,top=1.5cm,bottom=3cm]{geometry}
}

% Beamer version theme settings

%\useoutertheme[height=0pt,width=2cm,right]{sidebar}
%\usecolortheme{rose,sidebartab}
%\useinnertheme{circles}
%\usefonttheme[only large]{structurebold}

\setbeamercolor{sidebar right}{bg=black!15}
\setbeamercolor{structure}{fg=blue}
\setbeamercolor{author}{parent=structure}

\setbeamerfont{title}{series=\normalfont,size=\LARGE}
\setbeamerfont{title in sidebar}{series=\bfseries}
\setbeamerfont{author in sidebar}{series=\bfseries}
\setbeamerfont*{item}{series=}
\setbeamerfont{frametitle}{size=}
\setbeamerfont{block title}{size=\small}
\setbeamerfont{subtitle}{size=\normalsize,series=\normalfont}

\defbeamertemplate*{footline}{infolines theme}
 {
   \leavevmode%
   \hbox{%
   \begin{beamercolorbox}[wd=.333333\paperwidth,ht=2.25ex,dp=1ex,center]{author in head/foot}%
   %  \usebeamerfont{author in head/foot}\insertshortauthor~~\beamer@ifempty{\insertshortinstitute}{}{(\insertshortinstitute)}
   \end{beamercolorbox}%
   \begin{beamercolorbox}[wd=.333333\paperwidth,ht=2.25ex,dp=1ex,center]{title in head/foot}%
    % \usebeamerfont{title in head/foot}\insertshorttitle
   \end{beamercolorbox}%
   \begin{beamercolorbox}[wd=.333333\paperwidth,ht=2.25ex,dp=1ex,right]{date in head/foot}%
     %\usebeamerfont{date in head/foot}\insertshortdate{}\hspace*{2em}
     \insertshortlecture.\insertframenumber{} / \insertshortlecture.\inserttotalframenumber\hspace*{2ex} 
   \end{beamercolorbox}}%
   \vskip0pt%
 }
  


\setbeamertemplate{sidebar right}
{
  {\usebeamerfont{title in sidebar}%
    \vskip1.5em%
    \hskip3pt%
    \usebeamercolor[fg]{title in sidebar}%
    \insertshorttitle[width=2cm-6pt,center,respectlinebreaks]\par%
    \vskip1.25em%
  }%
  {%
    \hskip3pt%
    \usebeamercolor[fg]{author in sidebar}%
    \usebeamerfont{author in sidebar}%
    \insertshortauthor[width=2cm-2pt,center,respectlinebreaks]\par%
    \vskip1.25em%
  }%
  \hbox to2cm{\hss\insertlogo\hss}
  \vskip1.25em%
  \insertverticalnavigation{2cm}%
  \vfill
  \hbox to 2cm{\hfill\usebeamerfont{subsection in
      sidebar}\strut\usebeamercolor[fg]{subsection in
      sidebar}\insertshortlecture.\insertframenumber\hskip5pt}%
  \vskip3pt%
}%

\setbeamertemplate{title page}
{
  \vbox{}
  \vskip1em
  %{\huge Kapitel \insertshortlecture\par}
  {\usebeamercolor[fg]{title}\usebeamerfont{title}\inserttitle\par}%
  \ifx\insertsubtitle\@empty%
  \else%
    \vskip0.25em%
    {\usebeamerfont{subtitle}\usebeamercolor[fg]{subtitle}\insertsubtitle\par}%
  \fi%     
  \vskip1em\par
  %Vorlesung \emph{\lecturename}\ vom 
  \insertdate\par
  \vskip0pt plus1filll
  \leftskip=0pt plus1fill\insertauthor\par
  \insertinstitute\vskip1em
}

%\logo{\includegraphics[width=2cm]{beamerexample-lecture-logo.pdf}}



% Article version layout settings

\mode<article>

\makeatletter
\def\@listI{\leftmargin\leftmargini
  \parsep 0pt
  \topsep 5\p@   \@plus3\p@ \@minus5\p@
  \itemsep0pt}
\let\@listi=\@listI


\setbeamertemplate{frametitle}{\paragraph*{\insertframetitle\
    \ \small\insertframesubtitle}\ \par
}
\setbeamertemplate{frame end}{%
  \marginpar{\scriptsize\hbox to 1cm{\sffamily%
      \hfill\strut\insertshortlecture.\insertframenumber}\hrule height .2pt}}
\setlength{\marginparwidth}{1cm}
\setlength{\marginparsep}{1.5cm}

\def\@maketitle{\makechapter}

\def\makechapter{
  \newpage
  \null
  \vskip 2em%
  {%
    \parindent=0pt
    \raggedright
    \sffamily
    \vskip8pt
    %{\fontsize{36pt}{36pt}\selectfont Kapitel \insertshortlecture \par\vskip2pt}
    {\fontsize{24pt}{28pt}\selectfont \color{blue!50!black} \insertlecture\par\vskip4pt}
    {\Large\selectfont \color{blue!50!black} \insertsubtitle, \@date\par}
    \vskip10pt

    \normalsize\selectfont \@author\par\vskip1.5em
    %\hfill BLABLA
  }
  \par
  \vskip 1.5em%
}

\let\origstartsection=\@startsection
\def\@startsection#1#2#3#4#5#6{%
  \origstartsection{#1}{#2}{#3}{#4}{#5}{#6\normalfont\sffamily\color{blue!50!black}\selectfont}}

\makeatother

\mode
<all>




% Typesetting Listings

\usepackage{listings}
\lstset{language=Java}

\alt<presentation>
{\lstset{%
  basicstyle=\footnotesize\ttfamily,
  commentstyle=\slshape\color{green!50!black},
  keywordstyle=\bfseries\color{blue!50!black},
  identifierstyle=\color{blue},
  stringstyle=\color{orange},
  escapechar=\#,
  emphstyle=\color{red}}
}
{
  \lstset{%
    basicstyle=\ttfamily,
    keywordstyle=\bfseries,
    commentstyle=\itshape,
    escapechar=\#,
    emphstyle=\bfseries\color{red}
  }
}



% Common theorem-like environments

\theoremstyle{definition}
\newtheorem{exercise}[theorem]{\translate{Exercise}}




% New useful definitions:

\newbox\mytempbox
\newdimen\mytempdimen

\newcommand\includegraphicscopyright[3][]{%
  \leavevmode\vbox{\vskip3pt\raggedright\setbox\mytempbox=\hbox{\includegraphics[#1]{#2}}%
    \mytempdimen=\wd\mytempbox\box\mytempbox\par\vskip1pt%
    \fontsize{3}{3.5}\selectfont{\color{black!25}{\vbox{\hsize=\mytempdimen#3}}}\vskip3pt%
}}

\newenvironment{colortabular}[1]{\medskip\rowcolors[]{1}{blue!20}{blue!10}\tabular{#1}\rowcolor{blue!40}}{\endtabular\medskip}

\def\equad{\leavevmode\hbox{}\quad}

\newenvironment{greencolortabular}[1]
{\medskip\rowcolors[]{1}{green!50!black!20}{green!50!black!10}%
  \tabular{#1}\rowcolor{green!50!black!40}}%
{\endtabular\medskip}



\lecture[1]{2. Markgildi og samfelldni}{lecture-text}
\date{29. ágúst 2015}

\newcommand{\C}{{\mathbb  C}}
\newcommand{\Z}{{\mathbb Z}}
\newcommand{\R}{{\mathbb  R}}
\newcommand{\N}{{\mathbb  N}}
\newcommand{\Q}{{\mathbb Q}}

\begin{document}
\setcounter{tocdepth}{2}
\tableofcontents

\section{Markgildi og samfelldni}

%\subsection*{Inngangur}
\textbf{Nauðsynleg undirstaða}
\begin{itemize}
  \item Jafna línu, P.2
  \item Jafna hrings, P.3
  \item Hliðrun og skölun grafs, P.3
  \item (Stranglega) minnkandi og (stranglega) vaxandi föll, 2.8
  \item Jafnstæð og oddstæð föll, P.4
  \item Margliður; deiling, þáttun og rætur, P.6
  \item Tölugildisfallið, P.1
  \item Þríhyrningsójafnan, P.1
  \item Formerkjafallið, $sgn(x)$, P.5
\end{itemize}

\subsection{Markgildi}
\subsubsection{Óformleg skilgreining á markgildi}  
%Gerum ráð fyrir að fall $f$ sé
%skilgreint á opnu bili umhverfis punktinn $a$, nema hvað hugsanlega er
%$f(a)$ ekki skilgreint.  
Segjum að  $f(x)$
\emph{stefni á tölu} $L$ \emph{þegar} $x$ \emph{stefnir á} $a$, og ritum
$\lim_{x\rightarrow a} f(x)=L$, ef við getum tryggt að  $f(x)$ sé \emph{eins nálægt}
$L$ og við viljum bara með því að velja $x$ \emph{nógu nálægt} $a$.  

\pause

\subsubsection{Skilgreining: Markgildi} Gerum ráð fyrir að fall $f$ sé
skilgreint á opnu bili umhverfis punktinn $a$, nema hvað hugsanlega er
$f(a)$ ekki skilgreint. \pause 
Við segjum að $f(x)$
\emph{stefni á tölu} $L$ \emph{þegar} $x$ \emph{stefnir á} $a$, og ritum
$\lim_{x\rightarrow a} f(x)=L$, ef eftirfarandi skilyrði er uppfyllt:

\emph{Fyrir sérhverja tölu} $\epsilon>0$ \emph{er til tala} $\delta>0$ 
\emph{þannig að um öll} $x$ \emph{þannig að}
$$
0<|x-a|<\delta,\quad \text{ þá er } \quad |f(x)-L|<\epsilon.
$$

Við segjum að talan $L$ sé \emph{markgildi} $f(x)$ þegar $x$ stefnir á $a$.

\pause

\subsubsection{Athugasemd}  Þegar athugað er hvort markgildið
$\lim_{x\rightarrow a} f(x)$ er til og hvert gildi þess er þá skiptir
ekki máli hvort $f(a)$ er skilgreint eða ekki.

\subsection{Markgildi frá hægri}
\subsubsection{Óformleg skilgreining}   
Gerum ráð fyrir að fall $f$ sé
skilgreint á opnu bili $(a,b)$.  Segjum að  $f(x)$
\emph{stefni á tölu} $L$ \emph{þegar} $x$ \emph{stefnir á} $a$ \emph{frá hægri}, og ritum
$\lim_{x\rightarrow a^+} f(x)=L$, ef við getum tryggt að  $f(x)$ sé 
\emph{eins nálægt}
$L$ og við viljum bara með því að velja $x>a$ \emph{nógu nálægt} $a$. 

\pause

\subsubsection{Skilgreining (Markgildi frá hægri)} Gerum ráð fyrir að fall $f$ sé
skilgreint á opnu bili $(a,b)$.  Við segjum að $f(x)$
\emph{stefni á tölu} $L$ \emph{þegar} $x$ \emph{stefnir á} $a$ \emph{frá hægri}, og ritum
$\lim_{x\rightarrow a^+} f(x)=L$, ef eftirfarandi skilyrði er uppfyllt.

\emph{Fyrir sérhverja tölu} $\epsilon>0$ \emph{er til tala} $\delta>0$ \emph{þannig 
að um öll} $x$ \emph{þannig að} 
$$
a<x<a+\delta,\quad \text{ þá er } \quad |f(x)-L|<\epsilon.
$$

\subsection{Markgildi frá vinstri}
\subsubsection{Óformleg skilgreining}   
Gerum ráð fyrir að fall $f$ sé
skilgreint á opnu bili $(b,a)$.  Segjum að  $f(x)$
\emph{stefni á tölu} $L$ \emph{þegar} $x$ \emph{stefnir á} $a$ \emph{frá vinstri}, og ritum
$\lim_{x\rightarrow a^-} f(x)=L$, ef við getum tryggt að  $f(x)$ sé 
\emph{eins nálægt}
$L$ og við viljum bara með því að velja $x<a$ \emph{nógu nálægt} $a$. 

\pause

\subsubsection{Skilgreining -- Markgildi frá vinstri} Gerum ráð fyrir að fall $f$ sé
skilgreint á opnu bili $(b,a)$.  Við segjum að $f(x)$
\emph{stefni á tölu} $L$ \emph{þegar} $x$ \emph{stefnir á} $a$ \emph{frá vinstri}, og ritum
$\lim_{x\rightarrow a^-} f(x)=L$, ef eftirfarandi skilyrði er uppfyllt.

\emph{Fyrir sérhverja tölu} $\epsilon>0$ \emph{er til tala} $\delta>0$ \emph{þannig
  að um öll} $x$ \emph{þannig að} 
$$
a-\delta<x<a,\quad \text{ þá er } \quad |f(x)-L|<\epsilon.
$$

\subsection{Reiknireglur fyrir markgildi}
\subsubsection{Setning} Gerum ráð fyrir að fall $f$ sé
skilgreint á opnu bili umhverfis punktinn $a$, nema hvað hugsanlega er
$f(a)$ ekki skilgreint.  Þá er 
$$\lim_{x\rightarrow a} f(x)=L$$
ef og aðeins ef 
$$\lim_{x\rightarrow a^-} f(x)=L=\lim_{x\rightarrow a^+} f(x).$$

\subsubsection{Setning}   Gerum ráð fyrir að
$\lim_{x\rightarrow a}f(x)=L$ og að   $\lim_{x\rightarrow a}g(x)=M$.
Þá gildir

\begin{enumerate}[(i)]
  \item $\lim_{x\rightarrow a}\Big(f(x)+g(x)\Big)=L+M$;
  \item $\lim_{x\rightarrow a}\Big(f(x)-g(x)\Big)=L-M$;
  \item $\lim_{x\rightarrow a}f(x)g(x)=LM$;
  \item $\lim_{x\rightarrow a}kf(x)=kL$, þar sem $k$ fasti;
  \item $\lim_{x\rightarrow a}f(x)/g(x)=L/M$, að því gefnu að $M\neq 0$;
  \item Gerum ráð fyrir að $m$ og $n$ séu 
heiltölur þannig að $f(x)^{m/n}$ sé
skilgreint fyrir öll $x$ á bili $(b,c)$ umhverfis $a$ (en ekki
endilega fyrir $x=a$) og að $L^{m/n}$ sé skilgreint.
Þá er $\lim_{x\rightarrow a}f(x)^{m/n}=L^{m/n}$.
  \item Ef til er bil $(b,c)$ sem inniheldur $a$ þannig að 
$f(x)\leq g(x)$ fyrir öll $x\in (b,c)$, nema kannski $x=a$, þá er
$\lim_{x\rightarrow a}f(x)=L\leq M=\lim_{x\rightarrow a}g(x)$.
\end{enumerate}

\subsubsection{Athugasemd -- Varúð}
Liður (i) í setningunni á undan segir að ef markgildin
$\lim_{x\to a} f(x)$ og $\lim_{x\to a} g(x)$ eru til þá
sé markgildið $\lim_{x\to a} (f(x)+g(x))$ einnig til.

\pause
En hún segir {\bf ekki} að ef 
$f$ og $g$ eru föll þannig að markgildið
$\lim_{x\to a} (f(x)+g(x))$ er til, að þá séu markgildin
$\lim_{x\to a} f(x)$ og $\lim_{x\to a} g(x)$ einnig til.

\subsubsection{Setning -- Klemmureglan}Gerum ráð fyrir að $f(x)\leq
g(x)\leq h(x)$ fyrir öll $x$ á bili $(b, c)$ sem
inniheldur $a$, nema kannski $x=a$.  Gerum enn fremur ráð fyrir að 
$$\lim_{x\rightarrow a}f(x)=\lim_{x\rightarrow a}h(x)=L.$$
Þá er $\lim_{x\rightarrow a}g(x)=L$.


\subsection{Algeng markgildi}
\subsubsection{Dæmi}
\begin{enumerate}[(i)] 
  \item $\lim_{x \to a} c = c$, $c$ fasti
\pause
  \item $\lim_{x \to a} x = a$
  \item $\lim_{x \to a} |x| = |a|$
\pause
  \item $\lim_{x \to 0} \frac{|x|}{x}$ er ekki til
\pause
  \item $\lim_{x \to 0^-} \frac{|x|}{x} = -1$
  \item $\lim_{x \to 0^+} \frac{|x|}{x} = 1$
\end{enumerate}

\subsubsection{Sýnidæmi -- Markgildi með $\sin$}
\begin{enumerate}[(i)] 
  \item $$\lim_{x\to 0} \sin\left(\frac 1x\right) \quad \text{er ekki til}$$
  \pause
  \item $$\lim_{x\to 0} x\sin\left(\frac 1x\right) = 0$$
  \pause
  \item $$\lim_{x \to 0} \frac{\sin(x)}{x} = 1$$
\end{enumerate}

\subsubsection{Sýnidæmi -- Markgildi með $|x|$}
\begin{enumerate}[(i)] 
  \item $$\lim_{x\to 0^+} \frac x{|x|} = 1$$
  \pause
  \item $$\lim_{x\to 0^-} \frac x{|x|} = -1$$
  \pause
  \item $$\lim_{x\to 0} \frac x{|x|} \quad \text{er ekki til}$$
\end{enumerate}

\subsection{Markgildi þegar $x$ stefnir á $\infty$ og $-\infty$}
\begin{center}
  \includegraphics[width=8cm,keepaspectratio=true]{./myndir/kafli02/06_liminf.png}
  % Mynd_vorpunar.eps: 277x122 pixel, 300dpi, 2.35x1.03 cm, bb=0 0 277 122
\end{center}

\subsubsection{Óformleg skilgreining}
Gerum ráð fyrir að fall $f$ sé
skilgreint á bili $(a, \infty)$.  Segjum að  $f(x)$
\emph{stefni á tölu} $L$ \emph{þegar} $x$ \emph{stefnir á} $\infty$, og ritum
$\lim_{x\rightarrow \infty} f(x)=L$, ef við getum tryggt að  $f(x)$ sé eins
\emph{nálægt}
$L$ og við viljum bara með því að velja $x$ \emph{nógu stórt}.

\subsubsection{Skilgreining}
Gerum ráð fyrir að fall $f$ sé
skilgreint á bili $(a,\infty)$.  Við segjum að $f(x)$
\emph{stefni á tölu} $L$ \emph{þegar} $x$ \emph{stefnir á} $\infty$, og ritum
$\lim_{x\rightarrow \infty} f(x)=L$, ef eftirfarandi skilyrði er uppfyllt:

\emph{Fyrir sérhverja tölu} $\epsilon>0$ \emph{er til tala} $R$ \emph{þannig að um öll} $x>R$
\emph{gildir að} $|f(x)-L|<\epsilon$.

Fyrir $-\infty$ er þetta gert með sama sniði.

\subsubsection{Óformlega skilgreining}
Gerum ráð fyrir að fall $f$ sé
skilgreint á bili $(-\infty, a)$.  Segjum að  $f(x)$
\emph{stefni á tölu} $L$ \emph{þegar} $x$ \emph{stefnir á} $-\infty$, og ritum
$\lim_{x\rightarrow -\infty} f(x)=L$, ef við getum tryggt að  $f(x)$ sé eins
\emph{nálægt} $L$ og við viljum bara með því að velja $x$ sem \emph{nógu stóra} mínus tölu.

\pause
 
\subsubsection{Skilgreining}
Gerum ráð fyrir að fall $f$ sé
skilgreint á bili $(-\infty,a)$.  Við segjum að $f(x)$
\emph{stefni á tölu} $L$ \emph{þegar} $x$ \emph{stefnir á} $-\infty$, og ritum
$\lim_{x\rightarrow -\infty} f(x)=L$, ef eftirfarandi skilyrði er uppfyllt:

\emph{Fyrir sérhverja tölu} $\epsilon>0$ \emph{er til tala} $R$ \emph{þannig að um öll} $x<R$ 
\emph{gildir að}  $|f(x)-L|<\epsilon$.

\subsection{$\infty$ og $-\infty$ sem markgildi}
\subsubsection{$\infty$ sem markgildi -- Óformleg skilgreining}
Gerum ráð fyrir að fall $f$ sé
skilgreint á opnu bili umhverfis punktinn $a$, nema hvað hugsanlega er
$f(a)$ ekki skilgreint.  Segjum að  $f(x)$
\emph{stefni á} $\infty$ \emph{þegar} $x$ \emph{stefnir á} $a$, og ritum
$\lim_{x\rightarrow a} f(x)=\infty$, ef við getum tryggt að  $f(x)$ sé 
\emph{hversu stórt sem við viljum}  bara með því að velja $x$ \emph{nógu nálægt} $a$.  
 
\pause
\subsubsection{Skilgreining}
Gerum ráð fyrir að fall $f$ sé
skilgreint á opnu bili umhverfis punktinn $a$, nema hvað hugsanlega er
$f(a)$ ekki skilgreint.  Við segjum að $f(x)$
\emph{stefni á} $\infty$ \emph{þegar} $x$ \emph{stefnir á} $a$, og ritum
$\lim_{x\rightarrow a} f(x)=\infty$, ef eftirfarandi skilyrði er uppfyllt

\emph{fyrir sérhverja tölu} $B$ \emph{er til tala} $\delta>0$ \emph{þannig
  að um öll} $x$ \emph{þannig að}
$0<|x-a|<\delta$ 
\emph{gildir að} $f(x)>B$.


\subsubsection{Athugasemd -- Málvenja}  Athugið að $\infty$ er {\bf ekki} tala.  Þó 
að  $\lim_{x\rightarrow a} f(x)=\infty$ þá er samt sagt að
markgildið $\lim_{x\rightarrow a} f(x)$ sé ekki til.

\subsubsection{$-\infty$ sem markgildi -- Óformleg skilgreining}  Gerum ráð fyrir að fall $f$ sé
skilgreint á opnu bili umhverfis punktinn $a$, nema hvað hugsanlega er
$f(a)$ ekki skilgreint.  Segjum að  $f(x)$
\emph{stefni á} $-\infty$ \emph{þegar} $x$ \emph{stefnir á} $a$, og ritum
$\lim_{x\rightarrow a} f(x)=-\infty$, ef við getum tryggt að  $f(x)$ sé 
  \emph{hversu lítið sem við viljum}
 bara með því að velja $x$ \emph{nógu nálægt} $a$.  
 
\pause
\subsubsection{Skilgreining}
 Gerum ráð fyrir að fall $f$ sé
skilgreint á opnu bili umhverfis punktinn $a$, nema hvað hugsanlega er
$f(a)$ ekki skilgreint.  Við segjum að $f(x)$
\emph{stefni á} $-\infty$ \emph{þegar} $x$ \emph{stefnir á} $a$, og ritum
$\lim_{x\rightarrow a} f(x)=-\infty$, ef eftirfarandi skilyrði er uppfyllt

\emph{fyrir sérhverja tölu} $B$ \emph{er til tala} $\delta>0$ \emph{þannig
  að um öll} $x$ \emph{þannig að}

$0<|x-a|<\delta$ 
\emph{gildir að} $f(x)<B$.

\subsubsection{Athugasemd -- Málvenja}  Athugið að $-\infty$ er {\bf ekki} tala.  Þó 
að  $\lim_{x\rightarrow a} f(x)=-\infty$ þá er samt sagt að
markgildið $\lim_{x\rightarrow a} f(x)$ sé ekki til.
 

\subsection{Samfelldni}
\subsubsection{Skilgreining}
Látum $A\subseteq \R$ og $x\in A$.  Við segjum að $x$ sé \emph{innri
punktur} $A$ ef $A$ inniheldur opið bil umhverfis $x$, það er að
segja til er tala $\delta>0$ þannig að $(x-\delta, x+\delta)\subseteq
A$. 

\pause

Ef $x$ er ekki innri punktur $A$ og $x\in A$ þá segjum við að $x$ sé
\emph{jaðarpunktur} $A$.

\pause

\subsubsection{Skilgreining}
Látum $f$ vera fall og $c$ innri punkt skilgreiningarsvæðis $f$.  Sagt
er að $f$ sé \emph{samfellt í punktinum} $c$ ef
$$\lim_{x\rightarrow c}f(x)=f(c).$$

\subsubsection{Setning}
Látum $f$ og $g$ vera föll.  Gerum ráð fyrir að $c$ sé innri punktur
skilgreiningarsvæðis beggja fallanna og að bæði föllin séu samfelld í
punktinum $c$.  Þá eru eftirfarandi föll samfelld í $c$:
\begin{enumerate}[(i)] 
\pause
\item $f+g$ %(Fallið $f+g$ er skilgreint þannig að $(f+g)(x)=f(x)+g(x)$.)
\pause
\item $f-g$
\pause
\item $fg$
\pause
\item $kf$, þar sem $k$ er fasti
\pause
\item $f/g$, ef $g(c)\neq 0$
\pause
\item $\Big(f(x)\Big)^{1/n}$, að því gefnu að 
$f(c)>0$ ef $n$ er slétt tala og $f(c)\neq 0$ ef $n<0$.
\end{enumerate} 

\subsubsection{Setning -- Samskeyting samfelldra falla}
Látum $g$ vera fall sem er skilgreint á opnu bili umhverfis $c$ og
samfellt í $c$ og látum $f$ vera fall sem er skilgreint á opnu bili
umhverfis $g(c)$ og samfellt í $g(c)$.  Þá er fallið $f\circ g$
skilgreint á opnu bili umhverfis $c$ og er samfellt í $c$.  
%[Fallið $f\circ g$ er skilgreint þannig að $(f\circ g)(x)=f(g(x))$ fyrir þau
%gildi á $x$ þar sem bæði $g(x)$ og $f(g(x))$ eru skilgreind.]
% 
% Ef $f$ er samfellt í punkti $L$ og $\lim_{x\rightarrow c}g(x)=L$ þá er
% $$\lim_{x\rightarrow c}f\big(g(x)\big)=f(L)=f\bigg(\lim_{x\rightarrow c}g(x)\bigg).$$
% [Ekki forsenda varðandi þetta síðasta atriði að $g$ sé samfellt í $c$.]
 


\subsubsection{Athugasemd -- Hefð}
Ef fall er skilgreint með formúlu og skilgreingamengið er ekki
tilgreint sérstaklega, þá er venjan að líta alla þá punkta þar
sem formúlan gildir sem skilgreingarmengi fallsins
 
\subsubsection{Skilgreining}\label{skilgr:samfellt_fall}
 Við segjum að fall $f$ sé \emph{samfellt} ef það er samfellt í sérhverjum
punkti skilgreingarmengisins.

\subsubsection{Dæmi}
Eftirfarandi föll eru samfelld
\begin{enumerate}[(i)] 
\pause
  \item margliður
\pause
  \item ræð föll
\pause
  \item ræð veldi
\pause
  \item hornaföll; $\sin$, $\cos$, $\tan$
\pause
  \item tölugildisfallið $|x|$
\end{enumerate}

\subsubsection{Athugasemd -- Að búa til samfelld föll}
Með því að nota föllin úr dæminu á undan sem efnivið þá getum við búið til 
fjölda samfelldra fall með því að beita aðgerðunum úr Setningu 3.14 og 
Setningu 3.15.

\subsection{Hægri/vinstri samfelldni}	
Rifjum upp skilgreininguna á samfelldni.
\subsubsection{Skilgreining}
Látum $f$ vera fall og $c$ innri punkt skilgreiningarsvæðis $f$.  Sagt
er að $f$ sé \emph{samfellt í punktinum} $c$ ef
$$\lim_{x\rightarrow c}f(x)=f(c).$$
 
\pause
\subsubsection{Athugasemd}
Þessi skilgreining virkar aðeins fyrir innri punkta skilgreiningarsvæðisins.
Þannig að ef ætlunin er að rannsaka samfelldni í jaðarpunktum þá gengur
þessi skilgreining ekki. Hins vegar getum við útvíkkað skilgreininguna á samfelldni fyrir 
hægri og vinstri endapunkta bila með því að einskorða okkur
við markgildi frá vinstri og hægri.

\subsubsection{Skilgreining}
\begin{enumerate}[(i)]
\item Fall $f$ er \emph{samfellt frá hægri í punkti} $c$ ef
$\lim_{x\rightarrow c^+}f(x)=f(c)$.
  
Hér er gert ráð fyrir að fallið $f$ sé amk.~skilgreint á bilinu $[c, a)$.
\item Fall $f$ er \emph{samfellt frá vinstri í punkti} $c$ ef
$\lim_{x\rightarrow c^-}f(x)=f(c)$.
  
Hér er gert ráð fyrir að fallið $f$ sé amk.~skilgreint á bilinu $(a, c]$.
\end{enumerate}

Uppfærum nú skilgreiningu \ref{skilgr:samfellt_fall}.
 \subsubsection{Skilgreining (uppfærð)}
 Gerum ráð fyrir að $f$ sé fall sem er skilgreint á
 mengi $A$, þar sem $A$ er sammengi endanlega margra bila.
 Við segjum að fallið $f$ sé \emph{samfellt} ef það er samfellt í 
 öllum innri punktum skilgreingarmengisins, og ef það er samfellt
 frá hægri/vinstri í jaðarpunktum skilgreingarmengisins,
 eftir því sem við á.

\pause
\subsubsection{Athugasemd}
Ef fall er samfellt á opnu bili $(a,b)$, og ef $a<c<d<b$, þá er fallið einnig samfellt á bilinu $[c,d]$.

\subsection{Eiginleikar samfelldra falla}
\subsubsection{Setning -- Há- og lággildislögmálið}
Látum $f$ vera samfellt fall skilgreint á lokuðu takmörkuðu bili
$[a,b]$.  Þá eru til tölur $x_1$ og $x_2$ í $[a,b]$ þannig að 
fyrir allar tölur $x$ í $[a,b]$ er $$f(x_1)\leq f(x)\leq f(x_2).$$

Þetta þýðir að samfellt fall $f$ á lokuðu og takmörkuðu
bili $[a,b]$ tekur bæði hæsta og lægsta gildi á bilinu. \pause
Hæsta
gildið er þá $f(x_2)$ og lægsta gildið er $f(x_1)$.
 
\subsubsection{Athugasemd}
Það er mögulegt að fallið taki há/lággildi sitt í fleiri en einum punkti.

\subsubsection{Setning -- Milligildissetningin}
Látum $f$ vera samfellt fall skilgreint á lokuðu takmörkuðu bili
$[a,b]$.  Gerum ráð fyrir að $s$ sé tala sem liggur á milli $f(a)$ og
$f(b)$.  Þá er til tala $c$ sem liggur á milli $a$ og $b$ þannig að
$f(c)=s$. 
 
<iframe scrolling="no" src="https://tube.geogebra.org/material/iframe/id/zEQQcGcQ/width/1075/height/767/border/888888/rc/false/ai/false/sdz/true/smb/false/stb/false/stbh/true/ld/false/sri/true/at/auto" width="1075px" height="767px" style="border:0px;"> </iframe>>
\begin{center}
 \includegraphics[width=8cm,keepaspectratio=true]{./myndir/kafli02/10_Milligildissetn.png}
 % Mynd_vorpunar-eps-converted-to.pdf: 277x122 pixel, 300dpi, 2.35x1.03 cm, bb=0 0 277 122
\end{center}

\subsubsection{Fylgisetning}
Ef $P(x)=a_nx^n+a_{n-1}x^{n-1}+\cdots+a_1x+a_0$ er margliða af
oddatölu stigi, þá er til rauntala $c$ þannig að $P(c)=0$.

\subsubsection{Sönnun}
Gerum ráð fyrir að $a_n>0$.
Þá er $\lim_{x\to -\infty} P(x) = -\infty$ og
$\lim_{x\to \infty} P(x) = \infty$.
Það þýðir að til eru tölur $a$ og $b$ 
þannig að $P(a)<0$ og $P(b)>0$.
Með því að beita Milligildissetningunni á fallið $P$ á 
bilinu $[a,b]$ og með $s=0$ þá fæst að til er núllstöð 
á bilinu $[a,b]$.
	
Ef $a_n < 0$ þá víxlast markgildin að ofan en röksemdafærslan er
að öðru leyti eins.

\end{document}
